% Ejercicio 12, inciso b) — issue #7. Mismo contenido que inciso-b.md.
% Fragmento sin preámbulo: lo incluye latex/main.tex, que ya pone el título del inciso.

\begin{quote}
\itshape
Implementa con 4 tipos de elementos (CuentaAhorro, CuentaCorriente, CuentaInversión,
CuentaEmpresarial) y 3 operaciones (EjecutarTransacción, RevertirTransacción,
AuditarTransacción). Cada operación debe comportarse de manera diferente según el tipo de
cuenta.
\end{quote}

La implementación sigue el diseño del inciso~a: tiene las mismas clases y relaciones que sus
dos diagramas de clases (figuras~\ref{fig:e12-visitor} y~\ref{fig:e12-transacciones}). El
código completo está en la carpeta \texttt{codigo} del ejercicio, en el repositorio; aquí se
muestran solo los fragmentos relevantes.

\subsubsection{Cómo ejecutarlo}

Requisito: SDK de .NET~10. Desde \texttt{ejercicio-12-transacciones-bancarias/codigo/}:

\begin{lstlisting}[language={}, numbers=none, basicstyle=\ttfamily\footnotesize, columns=fullflexible, keepspaces=true]
dotnet run --project Banco.Procesamiento   # la simulación
dotnet test                                 # las pruebas unitarias (xUnit)
\end{lstlisting}

Hay que indicar \texttt{-{}-project} porque la carpeta tiene dos proyectos (el programa y
sus pruebas). Las fechas de la simulación son fijas, así que la salida es la misma en cada
ejecución.

\subsubsection{Estructura del código}

\begin{table}[H]
\centering
\small
\begin{tabularx}{\textwidth}{>{\raggedright\arraybackslash}p{6.4cm}
                             >{\raggedright\arraybackslash}X}
\toprule
\textbf{Carpeta / fichero} & \textbf{Clase del diagrama y rol} \\
\midrule
\texttt{Cuentas/ICuenta.cs} & \texttt{ICuenta}: Element \\
\texttt{Cuentas/Cuenta.cs} & \texttt{Cuenta}: base común de las cuentas (número, titular,
                             saldo e historial) \\
\texttt{Cuentas/Cuenta*.cs} & Las 4 cuentas: ConcreteElement \\
\texttt{Cuentas/Asiento.cs} & \texttt{Asiento}: lo que una transacción le hizo a una cuenta;
                              lo que se guarda para revertir y auditar \\
\addlinespace
\texttt{Visitantes/IVisitanteCuenta.cs} & \texttt{IVisitanteCuenta}: Visitor \\
\texttt{Visitantes/EjecutarTransaccion.cs}, \texttt{RevertirTransaccion.cs},
\texttt{AuditarTransaccion.cs} & Las 3 operaciones: ConcreteVisitor \\
\texttt{Visitantes/CalcularIntereses.cs} & ConcreteVisitor nuevo: demuestra la
                                           extensibilidad \\
\texttt{Visitantes/RegistroAuditoria.cs}, \texttt{EstadoTransaccion.cs},
\texttt{OperacionRechazadaException.cs} & Resultados de auditar y de rechazar \\
\addlinespace
\texttt{Transacciones/Transaccion.cs} y sus 4 subclases & \texttt{Transaccion},
  \texttt{Deposito}, \texttt{Retiro}, \texttt{Transferencia}, \texttt{PagoServicio}:
  peticiones encapsuladas como objetos (al estilo Command) \\
\texttt{Transacciones/Movimiento.cs}, \texttt{Naturaleza.cs} & Los débitos y créditos de
                                                                 cada transacción \\
\addlinespace
\texttt{ProcesadorTransacciones.cs}, \texttt{ResultadoTransaccion.cs} & Client: crea los
  visitantes y garantiza «todo o nada» \\
\texttt{Dinero.cs} & Clase auxiliar interna: redondeo y formato de importes (no aparece en
                     los diagramas) \\
\texttt{Program.cs} & La simulación \\
\texttt{Banco.Procesamiento.Tests/} & Pruebas unitarias \\
\bottomrule
\end{tabularx}
\caption{Ficheros del proyecto y clase del diagrama que implementa cada uno.}
\end{table}

\subsubsection{Cada operación se comporta distinto según la cuenta}

Las reglas de negocio son inventadas, pero explícitas: cada una está documentada en el
\texttt{Visitar} que la aplica, y sus valores son parámetros del constructor de la cuenta.

\begin{table}[H]
\centering
\footnotesize
\begin{tabularx}{\textwidth}{>{\raggedright\arraybackslash}p{2.1cm}
                             >{\raggedright\arraybackslash}X
                             >{\raggedright\arraybackslash}X
                             >{\raggedright\arraybackslash}X}
\toprule
\textbf{Cuenta} & \textbf{\texttt{EjecutarTransaccion}} & \textbf{\texttt{RevertirTransaccion}}
  & \textbf{\texttt{AuditarTransaccion}} \\
\midrule
\textbf{Ahorro}
  & Admite 3 débitos por mes (retiros, transferencias salientes y pagos), según la fecha de
    la transacción: el cuarto del mismo mes se rechaza y un mes nuevo empieza de cero. Sin
    sobregiro: nunca deja el saldo en negativo.
  & Devuelve el monto y, si era un débito del mes en curso, descuenta 1 de
    \texttt{RetirosDelMes}.
  & Conforme si la transacción (o su reversión) no dejó el saldo en negativo. \\
\addlinespace
\textbf{Corriente}
  & Permite quedar en negativo hasta $-500$. Cada débito que deja el saldo en negativo cobra
    un cargo fijo de 5 y suma 1 a \texttt{VecesEnSobregiro}; el tope se comprueba con el cargo
    incluido.
  & Devuelve el monto y el cargo y, si el débito había dejado la cuenta en negativo,
    descuenta 1 de \texttt{VecesEnSobregiro}.
  & Comprueba que el cargo se cobrara solo si el débito dejó la cuenta en negativo, y que el
    saldo no bajara del tope. \\
\addlinespace
\textbf{Inversión}
  & Los créditos tienen un mínimo de 100. Un débito antes de la fecha de vencimiento paga
    una penalización del 2\,\%.
  & Devuelve el monto y la penalización y la descuenta de
    \texttt{PenalizacionesCobradas}.
  & Comprueba la penalización según la fecha, el depósito mínimo y que el saldo no quedara
    negativo. \\
\addlinespace
\textbf{Empresarial}
  & Cada débito paga una comisión del 0,5\,\%. Los débitos de más de 10\,000 necesitan la
    autorización de alguien distinto de quien pide la transacción.
  & Devuelve el monto y la comisión y, si era un monto alto, descuenta 1 de
    \texttt{OperacionesAutorizadas}.
  & Comprueba la comisión, que los montos altos vengan autorizados por otra persona y que el
    saldo no quedara negativo; anota quién autorizó. \\
\bottomrule
\end{tabularx}
\caption{Reglas de cada operación según el tipo de cuenta.}
\end{table}

En las cuatro cuentas se cumple además lo siguiente:
\begin{itemize}
  \item \textbf{Revertir restaura exactamente el estado anterior}, saldo y contadores, porque
        aplica la operación inversa del \texttt{Asiento} que dejó la ejecución (inciso~a). La
        única excepción es deliberada: si se revierte un débito de ahorro de un mes que ya
        terminó, el contador del mes en curso no cambia.
  \item \textbf{Simplificación de la cuenta de ahorro:} solo lleva el conteo del mes más
        reciente, así que se supone que las transacciones llegan en orden de fecha.
  \item \textbf{Auditar no modifica la cuenta.} Por cada cuenta genera un
        \texttt{RegistroAuditoria} con quién pidió la transacción (\texttt{Usuario}), qué se
        pidió (\texttt{Descripcion}), cuándo (\texttt{Fecha}) y el resultado
        (\texttt{Estado}: aplicada, revertida o no aplicada), además de si lo que la
        transacción le hizo a la cuenta cumple las reglas de su tipo.
  \item \textbf{El estado es por cuenta, no por transacción.} \texttt{Revertida} incluye la
        compensación automática del «todo o nada»: en una transferencia rechazada, la cuenta
        que alcanzó a aplicar el débito queda como \texttt{Revertida} (con el rastro de la
        anulación, como un asiento contable de anulación) y la otra como
        \texttt{NoAplicada}.
  \item \textbf{La auditoría juzga la transacción, no el estado actual de la cuenta.} Revisa
        lo que dice el asiento: el cargo cobrado, la autorización y el saldo en que la
        transacción (o su reversión) dejó la cuenta. Una transacción correcta sigue siendo
        conforme aunque después otras cambien la cuenta.
\end{itemize}

\begin{lstlisting}[basicstyle=\ttfamily\footnotesize, columns=fullflexible, keepspaces=true, caption={\texttt{EjecutarTransaccion} en una cuenta corriente.}]
public void Visitar(CuentaCorriente cuenta)
{
    Movimiento movimiento = MovimientoPendiente(cuenta);
    if (movimiento.Naturaleza == Naturaleza.Credito)
    {
        Aplicar(cuenta, movimiento, cargo: 0m);
        return;
    }

    bool quedaEnNegativo = cuenta.Saldo - movimiento.Monto < 0m;
    decimal cargo = quedaEnNegativo ? cuenta.CargoPorSobregiro : 0m;
    ExigirSaldo(cuenta, movimiento.Monto + cargo, saldoMinimo: -cuenta.LimiteSobregiro,
                $"supera el sobregiro permitido de {Dinero.Formato(cuenta.LimiteSobregiro)}");
    Aplicar(cuenta, movimiento, cargo);
    if (quedaEnNegativo)
    {
        cuenta.VecesEnSobregiro++;
    }
}
\end{lstlisting}

El método recibe una \texttt{CuentaCorriente} y no una cuenta genérica: lo eligió el doble
despacho. Por eso puede leer \texttt{LimiteSobregiro} y \texttt{CargoPorSobregiro} sin
preguntar por el tipo ni hacer un cast. Ningún visitante usa \texttt{is}, \texttt{switch} ni
casts por tipo de cuenta. Se puede comprobar con la siguiente orden, que excluye los
comentarios de documentación (\texttt{///}): están en español, y algunas versiones de
\texttt{grep} encuentran «as» dentro de «así».

\begin{lstlisting}[language={}, numbers=none, basicstyle=\ttfamily\footnotesize, columns=fullflexible, keepspaces=true]
grep -rnE "\b(is|switch|as)\b" Banco.Procesamiento/Visitantes/ --include='*.cs' |
  grep -v "///"   # no devuelve nada
\end{lstlisting}

Ningún \texttt{Visitar} es una copia de otro: cada uno aplica la regla de su cuenta. Las partes
comunes (aplicar el asiento y exigir saldo) son métodos privados del visitante
(\texttt{Aplicar}, \texttt{ExigirSaldo}).

\subsubsection{«Todo o nada» en una transacción}

Una transferencia afecta a dos cuentas. \texttt{ProcesadorTransacciones} le pasa el
visitante a cada una y, si una rechaza su movimiento, revierte las que ya lo aplicaron, en
orden inverso, con \texttt{RevertirTransaccion}:

\begin{lstlisting}[basicstyle=\ttfamily\footnotesize, columns=fullflexible, keepspaces=true, caption={\texttt{ProcesadorTransacciones.Ejecutar}: todo o nada.}]
public ResultadoTransaccion Ejecutar(Transaccion transaccion)
{
    EjecutarTransaccion ejecutar = new(transaccion);
    List<ICuenta> aplicadas = [];
    try
    {
        foreach (ICuenta cuenta in transaccion.Cuentas)
        {
            cuenta.Aceptar(ejecutar);
            aplicadas.Add(cuenta);
        }

        return ResultadoTransaccion.Aplicada(transaccion);
    }
    catch (OperacionRechazadaException rechazo)
    {
        Deshacer(transaccion, aplicadas);
        return ResultadoTransaccion.Rechazada(transaccion, rechazo.Message);
    }
    catch
    {
        Deshacer(transaccion, aplicadas);
        throw;
    }
}
\end{lstlisting}

\begin{itemize}
  \item Las reglas comprueban todo \textbf{antes} de cambiar la cuenta: una cuenta que
        rechaza su movimiento queda intacta, y solo hay que deshacer las anteriores.
  \item Si la transacción falla por otra causa (por ejemplo, porque ya se había ejecutado en
        una de sus cuentas), también se deshace lo aplicado, y después se propaga el error.
  \item Los constructores de las transacciones garantizan que haya como mucho un movimiento
        por cuenta, y que los montos no tengan fracciones de centavo.
  \item La restricción del enunciado (issue~\#10) parte de este mecanismo para garantizar la
        atomicidad cuando las transacciones corren en paralelo.
\end{itemize}

\subsubsection{Extensibilidad: una operación nueva sin tocar las cuentas}

\texttt{CalcularIntereses} es una cuarta operación. Para agregarla \textbf{solo se creó esa
clase}, en la carpeta \texttt{Visitantes}: no se modificó ninguna cuenta, ninguna
transacción ni ninguno de los otros visitantes. También se comporta distinto según la cuenta:

\begin{table}[H]
\centering
\small
\begin{tabular}{ll}
\toprule
\textbf{Cuenta} & \textbf{Interés del mes} \\
\midrule
Ahorro & 3\,\% anual sobre el saldo positivo \\
Corriente & No paga intereses; si está en sobregiro, cobra el 18\,\% anual sobre lo que debe \\
Inversión & 6\,\% anual sobre el saldo positivo \\
Empresarial & 1\,\% anual, solo sobre lo que pase de 20\,000 \\
\bottomrule
\end{tabular}
\caption{Regla de \texttt{CalcularIntereses} según la cuenta.}
\end{table}

\begin{lstlisting}[basicstyle=\ttfamily\footnotesize, columns=fullflexible, keepspaces=true]
CalcularIntereses intereses = new();
foreach (Cuenta cuenta in cuentas)
{
    cuenta.Aceptar(intereses);   // el mismo Aceptar de siempre: las cuentas no cambiaron
}
\end{lstlisting}

\subsubsection{Resultados de la simulación}

\texttt{Program.cs} crea una cuenta de cada tipo y ejecuta 12 transacciones de los 4 tipos.
Algunas están pensadas para que se rechacen, una por cada regla:

\begin{lstlisting}[language={}, numbers=none, basicstyle=\ttfamily\scriptsize, columns=fullflexible, keepspaces=true]
--- 1. EjecutarTransaccion ---
  Aplicada   Depósito de 250.00 en AHO-001
  Aplicada   Retiro de 100.00 de AHO-001
  Aplicada   Retiro de 50.00 de AHO-001
  Aplicada   Pago de Electricidad por 80.00 desde AHO-001
  Rechazada  Retiro de 20.00 de AHO-001
             motivo: AHO-001: alcanzó el límite de 3 retiros del mes
  Aplicada   Retiro de 600.00 de COR-001
  Rechazada  Retiro de 300.00 de COR-001
             motivo: COR-001: supera el sobregiro permitido de 500.00 (saldo -305.00, a debitar 305.00)
  Aplicada   Retiro de 1000.00 de INV-001
  Rechazada  Transferencia de 50.00 de COR-001 a INV-001
             motivo: INV-001: el depósito mínimo es 100.00
  Rechazada  Transferencia de 15000.00 de EMP-001 a INV-001
             motivo: EMP-001: los débitos de más de 10000.00 necesitan la autorización de otra persona
  Aplicada   Transferencia de 15000.00 de EMP-001 a INV-001
  Aplicada   Pago de Telefonía por 1200.00 desde EMP-001
\end{lstlisting}

\begin{itemize}
  \item \textbf{8 aplicadas y 4 rechazadas.} Cada rechazo corresponde a una regla distinta:
        el límite de retiros, el tope de sobregiro, el depósito mínimo y la autorización.
  \item \textbf{Todo o nada.} La transferencia de 50 de COR-001 a INV-001 llegó a aplicar el
        débito en COR-001, con su cargo por sobregiro. Cuando INV-001 rechazó el crédito, el
        procesador revirtió ese débito, y COR-001 siguió en $-305$, con 1 vez en sobregiro.
  \item \textbf{Revertir} el retiro de 600 deja COR-001 en 300 y con 0 veces en sobregiro: el
        mismo estado que antes del retiro.
  \item \textbf{Auditar} genera 15 registros, uno por transacción y cuenta. Las cuentas a las
        que no llegó una transacción aparecen como \texttt{NoAplicada}. COR-001 aparece como
        \texttt{Revertida} en el retiro de 600 y en la transferencia rechazada, porque en las
        dos aplicó el débito y después se deshizo. Cada registro indica el saldo en que la
        transacción, o su reversión, dejó la cuenta.
  \item \textbf{Intereses del mes:} 2.55 (ahorro), 0.00 (corriente, ya sin sobregiro), 169.90
        (inversión) y 11.43 (empresarial).
\end{itemize}

\subsubsection{Pruebas}

\texttt{dotnet test} ejecuta \textbf{46 pruebas xUnit}, todas correctas. Para cada
combinación de cuenta y operación ($4 \times 3$) hay al menos una prueba:

\begin{table}[H]
\centering
\footnotesize
\begin{tabularx}{\textwidth}{>{\raggedright\arraybackslash}p{1.9cm}
                             >{\raggedright\arraybackslash}X
                             >{\raggedright\arraybackslash}X
                             >{\raggedright\arraybackslash}X}
\toprule
\textbf{Cuenta} & \textbf{Ejecutar} & \textbf{Ejecutar y luego revertir} & \textbf{Auditar} \\
\midrule
Ahorro & retiro cuenta como débito del mes; depósito no; rechaza el cuarto débito del mes; un
         mes nuevo empieza de cero; rechaza sin saldo; permite dejar el saldo en cero
       & un retiro y un depósito dejan el estado original; revertir un retiro de un mes
         anterior no cambia el mes en curso
       & registra quién, qué, cuándo y resultado sin modificar la cuenta \\
\addlinespace
Corriente & sin cargo si no queda en negativo; cargo al quedar en negativo; llega justo al
            tope; rechaza si pasa del tope
          & devuelve monto y cargo y descuenta el sobregiro
          & informa el sobregiro y el cargo sin modificar la cuenta; juzga la transacción y no
            el estado actual \\
\addlinespace
Inversión & penalización antes del vencimiento; sin penalización desde el vencimiento;
            rechaza depósito menor que el mínimo; rechaza si la penalización deja sin saldo
          & devuelve monto y penalización
          & señala el retiro anticipado y su penalización; marca no conforme la reversión que
            deja la cuenta en negativo \\
\addlinespace
Empresarial & cobra comisión en débitos; no en créditos; rechaza monto alto sin autorización;
              rechaza que alguien se autorice a sí mismo; acepta monto alto autorizado; un
              monto igual al umbral no necesita autorización
            & devuelve monto y comisión y descuenta la autorización
            & registra la comisión y quién autorizó sin modificar la cuenta \\
\bottomrule
\end{tabularx}
\caption{Pruebas de cada combinación de cuenta y operación.}
\end{table}

Cada prueba de «ejecutar y luego revertir» compara el saldo \textbf{y} los contadores de la
cuenta con los de antes de ejecutar. Además:
\begin{itemize}
  \item \textbf{\texttt{ProcesadorTransaccionesTests}} (7 pruebas):
        \begin{itemize}
          \item una transferencia aplica el débito y el crédito;
          \item si el crédito se rechaza, se revierte el débito (todo o nada);
          \item si la transacción falla por otra causa, se deshace lo aplicado y se propaga
                el error;
          \item revertir no borra lo que otra transacción hizo después en la misma cuenta;
          \item revertir o ejecutar dos veces la misma transacción falla sin tocar la cuenta;
          \item un depósito ya gastado se revierte igual, y la auditoría lo marca como no
                conforme.
        \end{itemize}
  \item \textbf{\texttt{TransaccionTests}} (3 pruebas): como mucho un movimiento por cuenta,
        ninguna transferencia a la misma cuenta y ningún monto con fracciones de centavo.
  \item \textbf{\texttt{DobleDespachoTests}}: a través de \texttt{ICuenta}, \texttt{Aceptar}
        llama al \texttt{Visitar} de la clase concreta.
  \item \textbf{\texttt{CalcularInteresesTests}} (3 pruebas): la operación nueva aplica la
        regla de cada cuenta, no da interés sobre un saldo negativo y no modifica nada.
\end{itemize}
